\chapter{总结与展望}

\section{工作总结}

本文所实现的数字化运营平台已经在公司内部使用，主要面向汽车电子软件项目管理和公司运营管理场景。本文围绕该平台的部分功能模块，完成了需求分析、系统设计、编码实现和功能测试等工作。

本文第一章分析了数字化运营平台的开发背景和国内外研究现状，介绍了论文的组织结构。第二章对系统用到的Spring Boot、Vue、MySQL、Redis、ECharts等关键技术做了简要说明。第三章从部门经理、项目经理、项目成员和质量人员四个角色，对BU工时统计分析、项目基本信息管理、缺陷管理和重大问题管理四个模块进行了功能性需求分析，同时从性能、可靠性和易用性等方面梳理了非功能性需求。第四章描述了 B/S 系统的五层架构，将整个平台按业务维度划分为系统管理、公司运营管理和项目运营管理三大模块，并说明了本文负责的后两个模块中各功能点的职责。第五章对四个模块的详细设计做了阐述，包括类图、时序图和数据库E-R模型。第六章展示了各模块的实现效果，通过功能测试用例和非功能性测试对系统进行了验证，其中针对缺陷深分页查询采用了延迟关联方案，优化后查询时间稳定在百毫秒级。

本次毕业设计具体负责的工作有：公司运营管理模块下的BU工时统计分析，支持按组织和年月维度查询各部门在业务单元上的工时投入，通过饼图和趋势图展示工时占比与变化趋势，利用Redis缓存减少数据库查询压力；项目运营管理模块下的项目基本信息维护，实现了项目的增删改查、条件筛选、Excel导出和交付历史追溯，采用主表加历史表的双写方式保留修改记录；缺陷管理模块，对接了Trinity、Jira和BugClose三个外部平台，采用策略模式与模板方法模式对多数据源查询接口进行重构，实现缺陷数据的统一查询和健康度分析；重大问题管理模块，实现了问题的新增、编辑、状态自动判定、飞书卡片通知和统计分析功能。

\section{工作展望}
本文设计实现的数字化运营平台已基本满足企业项目管理的日常需求，但在实际应用过程中仍暴露出一些不足之处，有待进一步优化和完善。未来的改进方向主要包括以下几个方面：
\begin{enumerate}[label=(\arabic*)]
	\item
	当前系统基于Spring Boot 2.1和Vue2构建，框架版本比较老旧，部分核心依赖已停止维护更新，存在潜在的安全风险。后续可以将前后端框架升级至Spring Boot 3.x，以便于引入 Spring AI 来让系统集成大模型服务，在缺陷分析、风险评估和项目报告生成等场景中探索智能化辅助能力。并迁移至 Java21 长期支持版本，因为这个版本的JDK提出了虚拟线程，非常适用于Web应用程序这种I/O密集型场景，前端每次启动速度都比较慢，可以考虑迁移至 Vue 3 配合 Vite 构建工具，以获得更好的性能和更完善的生态支持。
	
	\item
	目前平台已具备公司运营管理和项目运营管理两大功能域，能够为企业提供基本的数字化支撑。但从企业整体运营管理的视角来看，尚未覆盖企业运营管理的全部层级，后续可以在这个平台上继续扩展部门运营管理子系统，对部门级的人力资源效能、项目承接能力和团队协同效率做更细化的分析；同时可以考虑开发运营管理子系统，把需求吞吐量、代码提交频率、构建成功率等研发过程指标纳入平台，帮助团队发现开发流程中的瓶颈；还可以朝战略运营管理方向拓展，将各业务域的数据做更高层次的汇总，为管理层提供经营决策方面的数据支持。
\end{enumerate}